\documentclass[12pt, oneside]{book}

\usepackage{../style/mypackages}
\usepackage{../style/macro}
\usepackage{../style/font}
\usepackage{../style/theorem}
\usepackage{../style/format}

\begin{document}

\frontmatter
\pagenumbering{Roman}
\tableofcontents

\mainmatter
\pagenumbering{arabic}

\chapter{Curves}


\section{Parametrized Curves}

\begin{definition}
    A \textbf{parametrized differentiable curve} is a differentiable map
    \[
        \alpha:I\rightarrow \mathbb{R}^3
    \]
    from an open interval $I=(a,b)\subset \mathbb{R}$ into $\mathbb{R}^3$.
    If $\alpha(t)=(x(t),y(t),z(t))$, differentiability means that $x,y,z$ are
    differentiable functions.
\end{definition}

\begin{definition}
    The vector
    \[
        \alpha'(t)=(x'(t),y'(t),z'(t))
    \]
    is called the \textbf{tangent vector}, or \textbf{velocity vector}, of
    $\alpha$ at $t$.  The image set $\alpha(I)\subset \mathbb{R}^3$ is called
    the \textbf{trace} of $\alpha$.
\end{definition}

\begin{example}
    The curve
    \[
        \alpha(t)=(a\cos t,a\sin t,bt), \qquad t\in \mathbb{R},
    \]
    has as its trace a helix of pitch $2\pi b$ on the cylinder
    $x^2+y^2=a^2$.
\end{example}


\begin{definition}
    A parametrized differentiable curve
    $\alpha:I\rightarrow \mathbb{R}^3$ is said to be \textbf{regular} if
    \[
        \alpha'(t)\neq 0
    \]
    for every $t\in I$.
\end{definition}

\begin{definition}
    Let $\alpha:I\rightarrow \mathbb{R}^3$ be a regular parametrized curve.
    The \textbf{arc length} of $\alpha$ from $t_0\in I$ to $t\in I$ is
    \[
        s(t)=\int_{t_0}^{t}|\alpha'(u)|\,du .
    \]
    If $|\alpha'(s)|=1$, the curve is said to be \textbf{parametrized by arc
        length}.
\end{definition}

\begin{remark}
    Since $ds/dt=|\alpha'(t)|>0$, the arc length function has a differentiable
    inverse locally.  Thus every regular curve may be reparametrized by arc
    length.
\end{remark}


\section{The Local Theory of Curves Parametrized by Arc Length}

\begin{definition}
    Let $\alpha:I\rightarrow \mathbb{R}^3$ be a curve parametrized by arc
    length $s$.  The number
    \[
        \kappa(s):=|\alpha''(s)|
    \]
    is called the \textbf{curvature} of $\alpha$ at $s$.
\end{definition}

\begin{definition}
    Let $\alpha:I\rightarrow \mathbb{R}^3$ be a curve parametrized by arc
    length, with $\alpha''(s)\neq 0$ for all $s\in I$.  Define
    \[
        \mathbf{t}(s):=\alpha'(s),\qquad
        \mathbf{n}(s):=\frac{\alpha''(s)}{\kappa(s)},\qquad
        \mathbf{b}(s):=\mathbf{t}(s)\times \mathbf{n}(s).
    \]
    These are called the \textbf{unit tangent vector}, \textbf{principal normal
        vector}, and \textbf{binormal vector}, respectively.
\end{definition}

\begin{definition}
    The \textbf{torsion} $\tau(s)$ is defined by
    \[
        \mathbf{b}'(s)=\tau(s)\mathbf{n}(s).
    \]
\end{definition}

\begin{definition}
    Let $\alpha$ be as above.
    \begin{enumerate}
        \item The plane determined by $\mathbf{t}(s)$ and $\mathbf{n}(s)$ is
              called the \textbf{osculating plane} at $s$.
        \item The plane determined by $\mathbf{n}(s)$ and $\mathbf{b}(s)$ is
              called the \textbf{normal plane} at $s$.
        \item The plane determined by $\mathbf{t}(s)$ and $\mathbf{b}(s)$ is
              called the \textbf{rectifying plane} at $s$.
    \end{enumerate}
\end{definition}

\begin{theorem}[Frenet Formulas]
    Let $\alpha:I\rightarrow \mathbb{R}^3$ be a curve parametrized by arc
    length and assume $\alpha''(s)\neq 0$ for all $s\in I$.  Then
    \[
        \begin{aligned}
            \mathbf{t}' & = \kappa \mathbf{n},                  \\
            \mathbf{n}' & = -\kappa \mathbf{t}-\tau \mathbf{b}, \\
            \mathbf{b}' & = \tau \mathbf{n}.
        \end{aligned}
    \]
    Equivalently,
    \[
        \begin{pmatrix}
            \mathbf{t}' \\
            \mathbf{n}' \\
            \mathbf{b}'
        \end{pmatrix}
        =
        \begin{pmatrix}
            0       & \kappa & 0     \\
            -\kappa & 0      & -\tau \\
            0       & \tau   & 0
        \end{pmatrix}
        \begin{pmatrix}
            \mathbf{t} \\
            \mathbf{n} \\
            \mathbf{b}
        \end{pmatrix}.
    \]
\end{theorem}

\begin{theorem}[Fundamental Theorem of the Local Theory of Curves]
    Given differentiable functions $\kappa(s)>0$ and $\tau(s)$ on an interval
    $I$, there exists a regular parametrized curve
    $\alpha:I\rightarrow \mathbb{R}^3$ such that $s$ is arc length,
    $\kappa(s)$ is its curvature, and $\tau(s)$ is its torsion.  Moreover, any
    other such curve differs from $\alpha$ by a rigid motion.
\end{theorem}

\section{The Local Canonical Form}

\begin{remark}
    This section expresses a curve near a point in the coordinates determined by
    its Frenet trihedron.  It gives a local Taylor expansion whose coefficients
    are written in terms of curvature and torsion.
\end{remark}

\section{Global Properties of Plane Curves}

\begin{remark}
    This section studies global restrictions on plane curves, including results
    such as the isoperimetric inequality and the four-vertex theorem.
\end{remark}

\chapter{Regular Surfaces}

\section{Introduction}

The central object of this chapter is a regular surface in $\mathbb{R}^3$.
The goal is to develop a workable local model for surfaces using
parametrizations, tangent planes, differentials, and the first fundamental form.
The guiding idea is that a surface is a set which can be described locally by
two independent parameters.

\section{Regular Surfaces; Inverse Images of Regular Values}

\begin{definition}
    A subset $S\subset \mathbb{R}^3$ is a \textbf{regular surface} if, for each
    $p\in S$, there exist an open set $V\subset \mathbb{R}^3$ containing $p$,
    an open set $U\subset \mathbb{R}^2$, and a map
    $\mathbf{x}:U\rightarrow V\cap S$ such that
    \begin{enumerate}[label=(\roman*)]
        \item $\mathbf{x}$ is differentiable;
        \item $\mathbf{x}$ is a homeomorphism;
        \item for each $q\in U$, the differential $d\mathbf{x}_q$ is injective.
    \end{enumerate}
    Such a map $\mathbf{x}$ is called a \textbf{parametrization} of $S$.
\end{definition}

\begin{proposition}
    If $f:U\rightarrow \mathbb{R}$ is a differentiable function on an open set
    $U\subset \mathbb{R}^2$, then the graph
    \[
        \{(x,y,f(x,y)):(x,y)\in U\}
    \]
    is a regular surface.
\end{proposition}

\begin{proposition}
    If $f:U\subset \mathbb{R}^3\rightarrow \mathbb{R}$ is differentiable and
    $a\in f(U)$ is a regular value of $f$, then $f^{-1}(a)$ is a regular
    surface in $\mathbb{R}^3$.
\end{proposition}

\begin{example}
    The unit sphere
    \[
        S^2=\{(x,y,z)\in \mathbb{R}^3:x^2+y^2+z^2=1\}
    \]
    is a regular surface, since it is the inverse image of the regular value
    $1$ of $f(x,y,z)=x^2+y^2+z^2$.
\end{example}

\begin{example}
    A torus with radii $a>r>0$ can be written as the inverse image of $r^2$
    by
    \[
        f(x,y,z)=z^2+\left(\sqrt{x^2+y^2}-a\right)^2 .
    \]
    A useful parametrization is
    \[
        \mathbf{x}(u,v)=((a+r\cos u)\cos v,(a+r\cos u)\sin v,r\sin u),
    \]
    where the parameters are taken in suitable open intervals.
\end{example}

\begin{proposition}
    Let $S\subset \mathbb{R}^3$ be a regular surface and let $p\in S$.  Then
    some neighborhood of $p$ in $S$ is the graph of a differentiable function
    of one of the following forms:
    \[
        z=f(x,y),\qquad y=g(x,z),\qquad x=h(y,z).
    \]
\end{proposition}

\begin{proof}
    Let $\mathbf{x}(u,v)=(x(u,v),y(u,v),z(u,v))$ be a parametrization at
    $p$, and set $q=\mathbf{x}^{-1}(p)$.  Since $d\mathbf{x}_q$ is injective,
    at least one of
    \[
        \frac{\partial(x,y)}{\partial(u,v)},\qquad
        \frac{\partial(y,z)}{\partial(u,v)},\qquad
        \frac{\partial(z,x)}{\partial(u,v)}
    \]
    is nonzero at $q$.  If, for instance,
    $\partial(x,y)/\partial(u,v)\neq 0$, the inverse function theorem applies
    to $(u,v)\mapsto (x(u,v),y(u,v))$.  Locally one may solve
    $(u,v)=(u(x,y),v(x,y))$, and the surface is given by
    $z=z(u(x,y),v(x,y))$.  The other cases are the same after changing the
    projected coordinates.
\end{proof}

\begin{proposition}
    Let $S$ be a regular surface and let
    $\mathbf{x}:U\subset \mathbb{R}^2\rightarrow \mathbb{R}^3$ satisfy
    $\mathbf{x}(U)\subset S$.  If $\mathbf{x}$ is differentiable, one-to-one,
    and $d\mathbf{x}_q$ is injective for all $q\in U$, then
    $\mathbf{x}^{-1}$ is continuous.  Hence $\mathbf{x}$ is a parametrization
    of its image.
\end{proposition}

\begin{remark}
    The cone $z=\sqrt{x^2+y^2}$ is not a regular surface at the origin.  If it
    were regular there, it would locally be the graph of a differentiable
    function over one of the coordinate planes; the only possible graph over the
    $xy$-plane is not differentiable at $(0,0)$.
\end{remark}

\section{Change of Parameters; Differentiable Functions on Surfaces}

\begin{proposition}[Change of Parameters]
    Let $\mathbf{x}:U\subset \mathbb{R}^2\rightarrow S$ and
    $\mathbf{y}:V\subset \mathbb{R}^2\rightarrow S$ be parametrizations of a
    regular surface $S$.  If
    \[
        W=\mathbf{x}(U)\cap \mathbf{y}(V)\neq \varnothing,
    \]
    then the change of coordinates
    \[
        h=\mathbf{x}^{-1}\circ \mathbf{y}:
        \mathbf{y}^{-1}(W)\rightarrow \mathbf{x}^{-1}(W)
    \]
    is a diffeomorphism.
\end{proposition}

\begin{proof}
    Continuity follows from the fact that parametrizations are
    homeomorphisms.  For differentiability, fix a point in the overlap and
    choose coordinates so that
    $\partial(x,y)/\partial(u,v)\neq 0$ for the parametrization
    $\mathbf{x}(u,v)=(x(u,v),y(u,v),z(u,v))$.  Extend $\mathbf{x}$ to
    \[
        F(u,v,t)=(x(u,v),y(u,v),z(u,v)+t).
    \]
    The differential of $F$ has nonzero determinant at the point, so $F$ has a
    local differentiable inverse.  Locally, $h=F^{-1}\circ \mathbf{y}$ followed
    by projection to the first two coordinates.  The same argument applied to
    $h^{-1}$ proves that $h$ is a diffeomorphism.
\end{proof}

\begin{definition}
    Let $V\subset S$ be open and let $f:V\rightarrow \mathbb{R}$.  We say
    that $f$ is \textbf{differentiable at} $p\in V$ if, for some
    parametrization $\mathbf{x}:U\rightarrow S$ with $p\in \mathbf{x}(U)$,
    the function
    \[
        f\circ \mathbf{x}:U\rightarrow \mathbb{R}
    \]
    is differentiable at $\mathbf{x}^{-1}(p)$.
\end{definition}

\begin{remark}
    The change of parameters proposition shows that the definition does not
    depend on the parametrization.  If two coordinate systems are related by a
    diffeomorphism, differentiability in one system is equivalent to
    differentiability in the other.
\end{remark}

\begin{definition}
    A map $\varphi:S_1\rightarrow S_2$ between regular surfaces is
    \textbf{differentiable at} $p\in S_1$ if, for parametrizations
    $\mathbf{x}:U\rightarrow S_1$ and $\mathbf{y}:V\rightarrow S_2$ with
    $p\in \mathbf{x}(U)$ and $\varphi(p)\in \mathbf{y}(V)$, the local
    expression
    \[
        \mathbf{y}^{-1}\circ \varphi\circ \mathbf{x}
    \]
    is differentiable near $\mathbf{x}^{-1}(p)$.
\end{definition}

\begin{definition}
    A differentiable map $\varphi:S_1\rightarrow S_2$ is a
    \textbf{diffeomorphism} if it is one-to-one, onto, and its inverse is also
    differentiable.  In this case $S_1$ and $S_2$ are \textbf{diffeomorphic}.
\end{definition}

\section{The Tangent Plane; The Differential of a Map}

\begin{definition}
    Let $S$ be a regular surface and $p\in S$.  The \textbf{tangent plane}
    $T_pS$ is the set of tangent vectors at $p$ of differentiable curves on
    $S$ passing through $p$.
\end{definition}

\begin{proposition}
    If $\mathbf{x}:U\rightarrow S$ is a parametrization and
    $p=\mathbf{x}(u_0,v_0)$, then
    \[
        T_pS=d\mathbf{x}_{(u_0,v_0)}(\mathbb{R}^2)
        =
        \operatorname{span}\{\mathbf{x}_u(u_0,v_0),
        \mathbf{x}_v(u_0,v_0)\}.
    \]
    In particular $\{\mathbf{x}_u,\mathbf{x}_v\}$ is a basis of $T_pS$.
\end{proposition}

\begin{definition}
    Let $\varphi:S_1\rightarrow S_2$ be differentiable and let $p\in S_1$.
    The \textbf{differential} of $\varphi$ at $p$ is the linear map
    \[
        d\varphi_p:T_pS_1\rightarrow T_{\varphi(p)}S_2
    \]
    defined as follows: if $w=\alpha'(0)$ for a curve
    $\alpha:(-\epsilon,\epsilon)\rightarrow S_1$ with $\alpha(0)=p$, then
    \[
        d\varphi_p(w)=(\varphi\circ \alpha)'(0).
    \]
\end{definition}

\begin{remark}
    If $\mathbf{x}$ and $\mathbf{y}$ are local parametrizations of $S_1$ and
    $S_2$, and
    \[
        \Phi=\mathbf{y}^{-1}\circ \varphi\circ \mathbf{x},
    \]
    then $d\varphi_p$ is represented in the bases
    $\{\mathbf{x}_u,\mathbf{x}_v\}$ and $\{\mathbf{y}_\xi,\mathbf{y}_\eta\}$
    by the Jacobian matrix of $\Phi$.
\end{remark}

\begin{definition}
    A differentiable map $\varphi:S_1\rightarrow S_2$ is a
    \textbf{local diffeomorphism at} $p$ if there are neighborhoods
    $U_1$ of $p$ and $U_2$ of $\varphi(p)$ such that
    $\varphi:U_1\rightarrow U_2$ is a diffeomorphism.
\end{definition}

\begin{proposition}
    If $\varphi:S_1\rightarrow S_2$ is differentiable and
    $d\varphi_p:T_pS_1\rightarrow T_{\varphi(p)}S_2$ is an isomorphism, then
    $\varphi$ is a local diffeomorphism at $p$.
\end{proposition}

\section{The First Fundamental Form; Area}

\begin{definition}
    Let $S$ be a regular surface in $\mathbb{R}^3$ and $p\in S$.  The natural
    inner product of $\mathbb{R}^3$ induces on $T_pS$ an inner product
    $\langle,\rangle_p$.  The corresponding quadratic form
    \[
        I_p(w)=\langle w,w\rangle_p=|w|^2,\qquad w\in T_pS,
    \]
    is called the \textbf{first fundamental form} of $S$ at $p$.
\end{definition}

\begin{remark}
    If $\mathbf{x}(u,v)$ is a parametrization and
    $\alpha(t)=\mathbf{x}(u(t),v(t))$, then
    \[
        \begin{aligned}
            I_p(\alpha'(0))
             & =
            \left\langle
            \mathbf{x}_u u' + \mathbf{x}_v v',
            \mathbf{x}_u u' + \mathbf{x}_v v'
            \right\rangle_p              \\
             & = E(u')^2+2Fu'v'+G(v')^2,
        \end{aligned}
    \]
    where
    \[
        E=\langle \mathbf{x}_u,\mathbf{x}_u\rangle,\qquad
        F=\langle \mathbf{x}_u,\mathbf{x}_v\rangle,\qquad
        G=\langle \mathbf{x}_v,\mathbf{x}_v\rangle .
    \]
\end{remark}

\begin{definition}
    Let $w_1=a\mathbf{x}_u+b\mathbf{x}_v$ and
    $w_2=c\mathbf{x}_u+d\mathbf{x}_v$.  Their inner product is
    \[
        \langle w_1,w_2\rangle
        =Eac+F(ad+bc)+Gbd .
    \]
    Thus the first fundamental form is locally represented by the matrix
    \[
        \begin{pmatrix}
            E & F \\
            F & G
        \end{pmatrix}.
    \]
\end{definition}

\begin{proposition}
    Let $\alpha(t)=\mathbf{x}(u(t),v(t))$ be a regular curve on $S$.  Its
    arc length between $t=a$ and $t=b$ is
    \[
        L(\alpha)=\int_a^b
        \sqrt{
            E(u')^2+2Fu'v'+G(v')^2
        }\,dt.
    \]
\end{proposition}

\begin{proposition}
    If two regular curves on $S$ meet at $p$ with tangent vectors
    $w_1,w_2\in T_pS$, then their angle $\theta$ is determined by
    \[
        \cos \theta=
        \frac{\langle w_1,w_2\rangle_p}
        {\sqrt{\langle w_1,w_1\rangle_p}\sqrt{\langle w_2,w_2\rangle_p}}.
    \]
    In particular, the coordinate curves of $\mathbf{x}$ are orthogonal exactly
    when $F=0$.
\end{proposition}

\begin{definition}
    A bounded region $R\subset S$ contained in a coordinate neighborhood
    $\mathbf{x}(U)$ has \textbf{area}
    \[
        A(R)=\iint_{\mathbf{x}^{-1}(R)}
        |\mathbf{x}_u\times \mathbf{x}_v|\,du\,dv
        =
        \iint_{\mathbf{x}^{-1}(R)}\sqrt{EG-F^2}\,du\,dv.
    \]
\end{definition}

\begin{example}
    For a graph $\mathbf{x}(u,v)=(u,v,f(u,v))$ one has
    \[
        E=1+f_u^2,\qquad F=f_uf_v,\qquad G=1+f_v^2,
    \]
    hence
    \[
        dA=\sqrt{1+f_u^2+f_v^2}\,du\,dv.
    \]
\end{example}

\begin{example}
    For the sphere of radius $r$ parametrized by
    \[
        \mathbf{x}(u,v)=(r\sin u\cos v,r\sin u\sin v,r\cos u),
    \]
    where $0<u<\pi$, one obtains
    \[
        E=r^2,\qquad F=0,\qquad G=r^2\sin^2 u,
    \]
    and therefore $dA=r^2\sin u\,du\,dv$.
\end{example}

\section{Orientation of Surfaces}

\begin{definition}
    A regular surface $S$ is \textbf{orientable} if it can be covered by
    coordinate neighborhoods in such a way that every change of coordinates
    between overlapping charts has positive Jacobian.  Such a choice of charts
    is an \textbf{orientation} of $S$.
\end{definition}

\begin{definition}
    For a parametrization $\mathbf{x}(u,v)$, the associated unit normal field is
    \[
        N=
        \frac{\mathbf{x}_u\times \mathbf{x}_v}
        {|\mathbf{x}_u\times \mathbf{x}_v|}.
    \]
\end{definition}

\begin{proposition}
    A regular surface $S\subset \mathbb{R}^3$ is orientable if and only if
    there exists a differentiable field of unit normal vectors
    $N:S\rightarrow \mathbb{R}^3$.
\end{proposition}

\begin{proof}
    If the coordinate changes have positive Jacobian, the local normals
    $\mathbf{x}_u\times \mathbf{x}_v/|\mathbf{x}_u\times \mathbf{x}_v|$
    agree on overlaps and define a global unit normal field.  Conversely, if a
    global unit normal field $N$ is given, choose the order of the parameters in
    each connected coordinate neighborhood so that the local normal equals
    $N$.  On an overlap, a negative Jacobian would reverse the local normal,
    contradicting the choice.
\end{proof}

\begin{example}
    Graphs and level surfaces $f^{-1}(a)$, where $a$ is a regular value, are
    orientable.  For a level surface one may take
    \[
        N=\frac{\nabla f}{|\nabla f|}.
    \]
\end{example}

\begin{example}
    The Mobius strip is nonorientable.  When a unit normal is transported once
    around its central circle, it returns with the opposite sign, so no global
    continuous choice of unit normal can exist.
\end{example}

\section{A Characterization of Compact Orientable Surfaces}

\begin{definition}
    A regular surface $S\subset \mathbb{R}^3$ is \textbf{compact} if it is
    compact as a subset of $\mathbb{R}^3$.
\end{definition}

\begin{definition}
    Let $S\subset \mathbb{R}^3$ be a compact connected regular surface.  A
    point $p\in \mathbb{R}^3\setminus S$ is said to lie in an
    \textbf{inside component} if it belongs to a bounded connected component
    of $\mathbb{R}^3\setminus S$.  The unbounded component is called the
    \textbf{outside}.
\end{definition}

\begin{theorem}[Jordan-Brouwer Separation for Regular Surfaces]
    A compact connected regular surface
    $S\subset \mathbb{R}^3$ separates $\mathbb{R}^3$ into exactly two
    connected components: one bounded and one unbounded.  The surface $S$ is
    the common boundary of both components.
\end{theorem}

\begin{remark}
    In this book the separation theorem is used as a topological input.  It
    matches the geometric intuition that a closed surface has an inside and an
    outside.
\end{remark}

\begin{theorem}
    Every compact connected regular surface in $\mathbb{R}^3$ is orientable.
\end{theorem}

\begin{proof}
    Let $\Omega$ be the bounded component of
    $\mathbb{R}^3\setminus S$.  Near each point of $S$, the surface is a graph,
    and one can choose the unit normal pointing away from $\Omega$.  These
    local outward normals agree on overlaps and therefore define a global unit
    normal field.  By the previous characterization, $S$ is orientable.
\end{proof}

\section{A Geometric Definition of Area}

\begin{remark}
    The area formula
    \[
        A(R)=\iint_D|\mathbf{x}_u\times \mathbf{x}_v|\,du\,dv
    \]
    is not merely a formal definition.  It is the limit of areas of small
    parallelograms approximating the image of small rectangles in the parameter
    domain.
\end{remark}

\begin{proposition}
    Let $R=\mathbf{x}(D)$ be a region contained in a coordinate neighborhood,
    where $D\subset \mathbb{R}^2$ is a bounded domain.  If $D$ is divided into
    small rectangles and each image patch is approximated by the tangent
    parallelogram generated by $\mathbf{x}_u\Delta u$ and
    $\mathbf{x}_v\Delta v$, then the limiting sum of parallelogram areas is
    \[
        \iint_D|\mathbf{x}_u\times \mathbf{x}_v|\,du\,dv.
    \]
\end{proposition}

\begin{remark}
    Since
    \[
        |\mathbf{x}_u\times \mathbf{x}_v|^2
        =
        |\mathbf{x}_u|^2|\mathbf{x}_v|^2
        -
        \langle \mathbf{x}_u,\mathbf{x}_v\rangle^2
        =
        EG-F^2,
    \]
    the geometric area element agrees with the expression obtained from the
    first fundamental form.
\end{remark}

\section*{Appendix: A Brief Review of Continuity and Differentiability}
\addcontentsline{toc}{section}{Appendix: A Brief Review of Continuity and Differentiability}

\subsection*{Continuity}

\begin{definition}
    A ball in $\mathbb{R}^n$ with center $p$ and radius $\epsilon>0$ is
    \[
        B_\epsilon(p)=\{q\in \mathbb{R}^n:|q-p|<\epsilon\}.
    \]
    A set $U\subset \mathbb{R}^n$ is \textbf{open} if every point of $U$ is
    contained in some ball contained in $U$.
\end{definition}

\begin{definition}
    A map $F:U\subset \mathbb{R}^n\rightarrow \mathbb{R}^m$ is
    \textbf{continuous at} $p\in U$ if for every $\epsilon>0$ there exists
    $\delta>0$ such that
    \[
        F(B_\delta(p))\subset B_\epsilon(F(p)).
    \]
\end{definition}

\begin{proposition}
    A map $F=(f_1,\ldots,f_m):U\subset \mathbb{R}^n\rightarrow
        \mathbb{R}^m$ is continuous if and only if each component function
    $f_i:U\rightarrow \mathbb{R}$ is continuous.
\end{proposition}

\begin{proposition}
    The composition of continuous maps is continuous.
\end{proposition}

\begin{definition}
    A continuous one-to-one map $F:A\subset \mathbb{R}^n\rightarrow
        \mathbb{R}^m$ is a \textbf{homeomorphism onto its image} if
    $F^{-1}:F(A)\rightarrow A$ is continuous.
\end{definition}

\begin{theorem}[Intermediate Value Theorem]
    If $f:[a,b]\rightarrow \mathbb{R}$ is continuous and
    $f(a)f(b)<0$, then there exists $c\in(a,b)$ such that $f(c)=0$.
\end{theorem}

\begin{theorem}[Extreme Value Theorem]
    A continuous real-valued function on a closed interval $[a,b]$ attains its
    maximum and minimum.
\end{theorem}

\begin{theorem}[Heine-Borel]
    Every open cover of a closed interval $[a,b]$ admits a finite subcover.
\end{theorem}

\subsection*{Differentiability in \texorpdfstring{$\mathbb{R}^n$}{Rn}}

\begin{definition}
    In this book, \textbf{differentiable} means having continuous partial
    derivatives of all orders.
\end{definition}

\begin{definition}
    A map $F=(f_1,\ldots,f_m):U\subset \mathbb{R}^n\rightarrow
        \mathbb{R}^m$ is differentiable if each component function $f_i$ is
    differentiable.
\end{definition}

\begin{definition}
    Let $F:U\subset \mathbb{R}^n\rightarrow \mathbb{R}^m$ be differentiable.
    The \textbf{differential} of $F$ at $p$ is the linear map
    \[
        dF_p:\mathbb{R}^n\rightarrow \mathbb{R}^m
    \]
    defined by
    \[
        dF_p(w)=(F\circ \alpha)'(0),
    \]
    where $\alpha(0)=p$ and $\alpha'(0)=w$.
\end{definition}

\begin{remark}
    In coordinates, if $F=(f_1,\ldots,f_m)$, then $dF_p$ is represented by the
    Jacobian matrix
    \[
        \left(\frac{\partial f_i}{\partial x_j}(p)\right).
    \]
\end{remark}

\begin{theorem}[Chain Rule]
    If $F:U\subset \mathbb{R}^n\rightarrow \mathbb{R}^m$ and
    $G:V\subset \mathbb{R}^m\rightarrow \mathbb{R}^k$ are differentiable and
    $F(U)\subset V$, then
    \[
        d(G\circ F)_p=dG_{F(p)}\circ dF_p .
    \]
\end{theorem}

\begin{theorem}[Inverse Function Theorem]
    Let $F:U\subset \mathbb{R}^n\rightarrow \mathbb{R}^n$ be differentiable.
    If $dF_p$ is an isomorphism, then there are neighborhoods $U_0$ of $p$
    and $V_0$ of $F(p)$ such that
    $F:U_0\rightarrow V_0$ is a diffeomorphism.
\end{theorem}

\begin{definition}
    A value $a\in \mathbb{R}^m$ is a \textbf{regular value} of
    $F:U\subset \mathbb{R}^n\rightarrow \mathbb{R}^m$ if, for every
    $p\in F^{-1}(a)$, the differential $dF_p$ is onto.
\end{definition}

\chapter{The Geometry of the Gauss Map}
\section{The Definition of the Gauss Map and Its Fundamental Properties}

\begin{definition}
    An \textbf{orientation} of a regular surface $S$ is a differentiable choice
    of a unit normal vector $N(p)$ at every $p\in S$.  A surface admitting such
    a choice is \textbf{orientable}.  For an oriented parametrization
    $x:U\rightarrow S$,
    \[
        N=\frac{x_u\wedge x_v}{|x_u\wedge x_v|}.
    \]
\end{definition}

Every regular surface is locally orientable, although it need not be globally
orientable; the M\"obius strip is the standard example.  An orientation of
$S$ determines an orientation of $T_pS$: a basis $\{v,w\}$ is positive when
$\langle v\wedge w,N(p)\rangle>0$.

\begin{definition}
    The \textbf{Gauss map} of the oriented surface $S$ is
    \[
        N:S\longrightarrow\mathbb{S}^2,
        \qquad p\longmapsto N(p).
    \]
    Its differential may be regarded as an endomorphism
    \[
        dN_p:T_pS\longrightarrow T_pS,
    \]
    since $T_{N(p)}\mathbb{S}^2$ and $T_pS$ are the same vector subspace of
    $\mathbb{R}^3$.
\end{definition}




\begin{proposition}
    The linear map $dN_p:T_pS\rightarrow T_pS$ defined by
    \begin{equation*}
        dN_p(x_u)=N_u,\qquad dN_p(x_v)=N_v
    \end{equation*}
    is self-adjoint with respect to the inner product of $\mathbb{R}^3$.
\end{proposition}

\begin{proof}
    In a parametrization $x(u,v)$, $dN_p(x_u)=N_u$ and
    $dN_p(x_v)=N_v$.  Differentiating
    $\langle N,x_u\rangle=\langle N,x_v\rangle=0$ gives
    \[
        \langle N_u,x_v\rangle=-\langle N,x_{uv}\rangle
        =\langle x_u,N_v\rangle.
    \]
    This proves the assertion on the basis $\{x_u,x_v\}$.
\end{proof}

\begin{definition}
    The \textbf{shape operator} (or \textbf{Weingarten map}) is
    \[
        W_p:=-dN_p:T_pS\longrightarrow T_pS.
    \]
    \begin{enumerate}

        \item
              The \textbf{Gaussian curvature} $K:= \det (W_p) $ and the mean curvature
              $H:=\frac12\operatorname{tr}(W_p)$ .
        \item
              The \textbf{principal curvatures} $k_1, k_2$ are the eigenvalues of $W_p$, and the corresponding orthonormal  eigenvectors are called \textbf{principal directions}.
        \item The \textbf{second fundamental form} is the quadratic form
              \[
                  \mathrm{II}_p(v)
                  :=
                  \langle W_p v,v\rangle
                  =-\langle dN_p(v),v\rangle.
              \]
              with representation in the basis $\{x_u,x_v\}$ given by the matrix
              \[
                  \begin{pmatrix}
                      e & f \\
                      f & g
                  \end{pmatrix},
              \]
              where $e=\langle x_{uu},N\rangle$, $f=\langle x_{uv},N\rangle$, and $g=\langle x_{vv},N\rangle$.
    \end{enumerate}
\end{definition}

\begin{definition}
    A point $p\in S$ is
    \begin{enumerate}
        \item \textbf{elliptic} if $K(p)>0$;
        \item \textbf{hyperbolic} if $K(p)<0$;
        \item \textbf{parabolic} if $K(p)=0$ and $dN_p\neq0$;
        \item \textbf{planar} if $dN_p=0$.
        \item \textbf{umbilical} if $k_1(p)=k_2(p)$.
    \end{enumerate}
\end{definition}

\section{Curve Theory on Surfaces}

\begin{definition}
    Let $C\subset S$ be a regular curve through $p$,, with curvature $k$ and principal normal $n$(or say curvature vecotr ).
    The \textbf{normal curvature} of $C$ in $S$ at $p$ is
    \[
        k_n:=k\langle n,N\rangle=\langle\alpha''(0),N(p)\rangle.
    \]
    \begin{remark}
        Changing the orientation of the curve does not change $k_n$, while changing
        the orientation of the surface changes its sign.
    \end{remark}
\end{definition}





The intersection of $S$ with the plane spanned by $v$ and $N(p)$ is the
\textbf{normal section} along $v$.  Its curvature at $p$ is $|k_n(v)|$.



\begin{proposition}[Olinde Rodrigues]
    A connected regular curve $\alpha(t)$ on $S$ is a line of curvature if
    and only if
    \[
        N'(t)=\lambda(t)\alpha'(t).
    \]
    In that case the principal curvature along $\alpha'(t)$ is
    $-\lambda(t)$.
\end{proposition}

If a unit vector $v$ makes an angle $\theta$ with $e_1$, then
\[
    v=e_1\cos\theta+e_2\sin\theta.
\]
Consequently one obtains Euler's formula
\begin{equation}
    k_n(v)=k_1\cos^2\theta+k_2\sin^2\theta.
    \label{eq:euler-normal-curvature}
\end{equation}



The value of $K$ is independent of the orientation of $S$, whereas the signs
of $H,k_1,k_2$ change when $N$ is replaced by $-N$.



Planes and spheres consist entirely of umbilical points.  On a cylinder every
point is parabolic; its principal curvatures are $0$ and the reciprocal of the
radius, up to the chosen orientation.

\begin{proposition}
    If every point of a connected regular surface $S$ is umbilical, then $S$
    is contained in a plane or in a sphere.
\end{proposition}

\begin{proof}
    Locally $dN=\lambda I$.  From $N_u=\lambda x_u$ and
    $N_v=\lambda x_v$, equality of mixed derivatives yields
    $\lambda_u=\lambda_v=0$.  If $\lambda=0$, $N$ is constant and the patch
    lies in a plane.  If $\lambda\neq0$, then
    $x-\lambda^{-1}N$ is constant, so the patch lies in a sphere of radius
    $|\lambda|^{-1}$.  Connectedness extends the conclusion to all of $S$.
\end{proof}

\begin{definition}
    A nonzero direction $v\in T_pS$ is an \textbf{asymptotic direction} if
    $\mathrm{II}_p(v)=0$.  A curve tangent everywhere to asymptotic directions
    is an \textbf{asymptotic curve}.
\end{definition}

There are no asymptotic directions at an elliptic point, exactly two at a
hyperbolic point, and one at a parabolic point.  At a planar point every
direction is asymptotic.

\begin{definition}
    The \textbf{Dupin indicatrix} at $p$ is the set of vectors $w\in T_pS$
    satisfying $\mathrm{II}_p(w)=\pm1$.  In principal coordinates
    $w=\xi e_1+\eta e_2$, its equation is
    \[
        k_1\xi^2+k_2\eta^2=\pm1.
    \]
\end{definition}

It is an ellipse at an elliptic point and a pair of hyperbolas with common
asymptotes at a hyperbolic point.  At a parabolic point it degenerates into a
pair of parallel lines.  These pictures encode the asymptotic directions and
the variation of normal curvature.

\begin{definition}
    Nonzero vectors $v,w\in T_pS$ are \textbf{conjugate} if
    \[
        \langle dN_p(v),w\rangle=0,
    \]
    equivalently $\langle A_pv,w\rangle=0$.
\end{definition}

Principal directions are conjugate; an asymptotic direction is conjugate to
itself.  If $v$ and $w$ make angles $\theta$ and $\varphi$ with $e_1$, then
they are conjugate precisely when
\[
    k_1\cos\theta\cos\varphi
    +k_2\sin\theta\sin\varphi=0.
\]

\section{The Gauss Map in Local Coordinates}

Let $x(u,v)$ be a positively oriented parametrization.  Write
\[
    E=\langle x_u,x_u\rangle,
    \qquad F=\langle x_u,x_v\rangle,
    \qquad G=\langle x_v,x_v\rangle
\]
for the coefficients of the first fundamental form.  The coefficients of the
second fundamental form are
\[
    e=\langle N,x_{uu}\rangle=-\langle N_u,x_u\rangle,
    \quad
    f=\langle N,x_{uv}\rangle=-\langle N_u,x_v\rangle,
    \quad
    g=\langle N,x_{vv}\rangle=-\langle N_v,x_v\rangle.
\]
Thus, for $w=x_u u'+x_v v'$,
\begin{equation}
    \mathrm{I}(w)=E(u')^2+2Fu'v'+G(v')^2,
    \qquad
    \mathrm{II}(w)=e(u')^2+2fu'v'+g(v')^2.
    \label{eq:fundamental-forms}
\end{equation}

Writing
\[
    N_u=a_{11}x_u+a_{21}x_v,
    \qquad
    N_v=a_{12}x_u+a_{22}x_v,
\]
the matrix of $dN$ in the basis $\{x_u,x_v\}$ is
\begin{equation}
    [dN]=\frac{1}{EG-F^2}
    \begin{pmatrix}
        fF-eG & gF-fG \\
        eF-fE & fF-gE
    \end{pmatrix}.
    \label{eq:weingarten-matrix}
\end{equation}
Equivalently, the Weingarten equations are obtained from
\[
    [A]=[\mathrm{I}]^{-1}[\mathrm{II}],
    \qquad A=-dN.
\]

It follows immediately that
\begin{equation}
    K=\frac{eg-f^2}{EG-F^2},
    \qquad
    H=\frac{eG-2fF+gE}{2(EG-F^2)}.
    \label{eq:KH-local}
\end{equation}
The principal curvatures are the roots of
\[
    k^2-2Hk+K=0,
    \qquad
    k_{1,2}=H\pm\sqrt{H^2-K}.
\]


\begin{proposition}
    The representation of $W_p$ in the basis $\{x_u,x_v\}$ is
    \begin{equation*}
        [W]=[\mathrm{I}]^{-1}[\mathrm{II}]
    \end{equation*}
\end{proposition}



\begin{proposition}
    Near an elliptic point the surface lies on one side of its tangent plane.
    Every neighborhood of a hyperbolic point contains points on both sides of
    its tangent plane.
\end{proposition}

\begin{proof}
    Choose $x(0,0)=p$ and let
    $d(u,v)=\langle x(u,v)-x(0,0),N(p)\rangle$.  Taylor's formula gives
    \[
        d(u,v)=\frac12(eu^2+2fuv+gv^2)+R,
        \qquad \frac{R}{u^2+v^2}\longrightarrow0.
    \]
    At an elliptic point the quadratic part has a fixed sign; at a hyperbolic
    point it takes both signs.
\end{proof}

The differential equation of the asymptotic curves is
\begin{equation}
    e(du)^2+2f\,du\,dv+g(dv)^2=0.
    \label{eq:asymptotic-curves}
\end{equation}
Near a hyperbolic point, the coordinate curves are asymptotic precisely when
$e=g=0$.

The differential equation of the lines of curvature is
\begin{equation}
    (fE-eF)(du)^2+(gE-eG)du\,dv+(gF-fG)(dv)^2=0.
    \label{eq:curvature-lines}
\end{equation}
Near a nonumbilical point, the coordinate curves are lines of curvature if and
only if $f=F=0$.

\begin{proposition}
    Let $K(p)\neq0$.  For sufficiently small regions $B\subset S$ containing
    $p$, let $A(B)$ be their area and let $A(N(B))$ be the signed area of their
    spherical images.  Then
    \[
        K(p)=\lim_{B\to p}\frac{A(N(B))}{A(B)}.
    \]
    The Gauss map preserves orientation where $K>0$ and reverses orientation
    where $K<0$.
\end{proposition}

\section{Vector Fields}

\begin{definition}
    A differentiable \textbf{vector field} on an open set
    $U\subset\mathbb{R}^2$ is a map
    \[
        w(x,y)=(a(x,y),b(x,y)).
    \]
    A curve $\alpha(t)$ is a \textbf{trajectory} of $w$ if
    $\alpha'(t)=w(\alpha(t))$, or equivalently
    \[
        \frac{dx}{dt}=a(x,y),
        \qquad
        \frac{dy}{dt}=b(x,y).
    \]
\end{definition}

\begin{theorem}[Local existence and uniqueness]
    Given $p\in U$, there is a trajectory $\alpha:I\rightarrow U$ with
    $\alpha(0)=p$.  Any two such trajectories agree wherever both are
    defined.
\end{theorem}

\begin{theorem}[Differentiable dependence]
    Near each $p\in U$ there are a neighborhood $V$, an interval $I$, and a
    differentiable map $\Phi:V\times I\rightarrow U$ such that
    \[
        \Phi(q,0)=q,
        \qquad
        \frac{\partial\Phi}{\partial t}(q,t)=w(\Phi(q,t)).
    \]
    The map $\Phi$ is the local flow of $w$.
\end{theorem}

\begin{lemma}
    If $w(p)\neq0$, there is a neighborhood $W$ of $p$ and a differentiable
    function $f:W\rightarrow\mathbb{R}$ such that $df\neq0$ and $f$ is
    constant on every trajectory of $w$ in $W$.
\end{lemma}

Such an $f$ is a \textbf{local first integral}.  It provides a coordinate
transverse to the trajectories.

\begin{definition}
    A differentiable \textbf{field of directions} assigns to each $p\in U$ a
    line through the origin of $T_pU$.  Locally it is represented by a
    nowhere-zero vector field.  A regular curve tangent to the assigned line
    at every point is an \textbf{integral curve}.
\end{definition}

The equation
\[
    a(x,y)\,dx+b(x,y)\,dy=0
\]
describes the direction spanned by $(b,-a)$.  The same definitions transfer
to a regular surface by using local coordinates.

\begin{definition}
    A vector field on an open set $U\subset S$ assigns
    $w(p)\in T_pS$ to every $p\in U$.  It is differentiable if in every local
    parametrization it can be written
    \[
        w=a(u,v)x_u+b(u,v)x_v
    \]
    with differentiable coefficients $a,b$.
\end{definition}

\begin{theorem}[Coordinates adapted to two fields]
    Let $w_1,w_2$ be differentiable vector fields on $U\subset S$ which are
    linearly independent at $p$.  There is a parametrization of a neighborhood
    of $p$ whose two families of coordinate curves are tangent to the
    directions determined by $w_1$ and $w_2$.
\end{theorem}

\begin{proof}
    Choose local first integrals $f_1,f_2$ of $w_1,w_2$.  The map
    $q\mapsto(f_1(q),f_2(q))$ has nonsingular differential at $p$ and hence is
    a local diffeomorphism.  Its inverse is the desired parametrization.
\end{proof}

Important consequences are the following.

\begin{corollary}
    Given two distinct differentiable fields of directions near $p$, there are
    local coordinates whose coordinate curves are their integral curves.
\end{corollary}

\begin{corollary}
    Every point of a regular surface has an orthogonal parametrization.
\end{corollary}

\begin{corollary}
    Near a hyperbolic point there is a parametrization whose coordinate curves
    are asymptotic curves.
\end{corollary}

\begin{corollary}
    Near a nonumbilical point there is a parametrization whose coordinate
    curves are lines of curvature.
\end{corollary}

\section{Ruled Surfaces and Minimal Surfaces}

\subsection{Ruled Surfaces}

\begin{definition}
    A one-parameter family of lines is described by a point $\alpha(t)$ and a
    nonzero direction $w(t)$.  The associated \textbf{ruled surface} is
    \[
        x(t,v)=\alpha(t)+v w(t).
    \]
    The lines $v\mapsto x(t,v)$ are the \textbf{rulings}, and $\alpha$ is a
    \textbf{directrix}.
\end{definition}

Cylinders, cones, tangent surfaces, the hyperbolic paraboloid, the helicoid,
and the one-sheeted hyperboloid are ruled surfaces.  Singular points are
allowed in this definition.

Assume $|w|=1$ and $w'\neq0$.  The ruled surface is then called
\textbf{noncylindrical}.  Its distinguished directrix
\begin{equation}
    \beta(t)=\alpha(t)-
    \frac{\langle\alpha'(t),w'(t)\rangle}{|w'(t)|^2}w(t)
    \label{eq:striction-line}
\end{equation}
is independent of the original choice of directrix and satisfies
$\langle\beta',w'\rangle=0$.  It is the \textbf{line of striction}, and its
points are the central points of the rulings.

Writing $x(t,u)=\beta(t)+uw(t)$, define the \textbf{distribution parameter}
\[
    \lambda(t)=\frac{(\beta'(t),w(t),w'(t))}{|w'(t)|^2},
\]
where $(a,b,c)=\langle a\wedge b,c\rangle$.  Then
\[
    |x_t\wedge x_u|^2=(\lambda^2+u^2)|w'|^2.
\]
Thus singular points, if any, lie on the line of striction and occur where
$\lambda=0$.

At every regular point,
\begin{equation}
    K(t,u)=-\frac{\lambda(t)^2}{(\lambda(t)^2+u^2)^2}\leq0.
    \label{eq:ruled-curvature}
\end{equation}
When $\lambda\neq0$, $|K|$ is maximal at the central point of each ruling.
If $\theta$ is the angle between the normals at $x(t,u)$ and $x(t,0)$, then
\[
    \tan\theta=\frac{u}{|\lambda|}.
\]

\begin{definition}
    A ruled surface $x(t,v)=\alpha(t)+vw(t)$ is \textbf{developable} if
    \[
        (w,w',\alpha')=0.
    \]
\end{definition}

At its regular points a developable surface has $K=0$.  Away from exceptional
accumulations of singularities, a developable surface is locally a piece of a
cylinder, a cone, or a tangent surface.  For a noncylindrical developable,
$\lambda=0$, so its line of striction is its singular locus.

\subsection{Minimal Surfaces}

\begin{definition}
    A regular parametrized surface is \textbf{minimal} if its mean curvature
    vanishes identically.
\end{definition}

Let $D$ be a bounded domain in the parameter plane and let $h$ be a
differentiable function.  A normal variation is
\[
    x_t(u,v)=x(u,v)+t h(u,v)N(u,v).
\]
If $A(t)$ is the area of $x_t(D)$, the first variation of area is
\begin{equation}
    A'(0)=-2\iint_D hH\sqrt{EG-F^2}\,du\,dv.
    \label{eq:first-variation-area}
\end{equation}

\begin{proposition}
    A regular parametrized surface is minimal if and only if
    $A'(0)=0$ for every bounded domain $D$ and every normal variation
    supported in $D$.
\end{proposition}

Thus minimal surfaces are critical points of area, although a critical point
need not always be an area minimum.  The \textbf{mean curvature vector}
$\mathbf H=HN$ points in a direction in which area initially decreases.

\begin{definition}
    A parametrization $x(u,v)$ is \textbf{isothermal} if
    \[
        E=G=\lambda^2,
        \qquad F=0.
    \]
\end{definition}

\begin{proposition}
    For an isothermal parametrization,
    \[
        x_{uu}+x_{vv}=2\lambda^2\mathbf H.
    \]
\end{proposition}

\begin{corollary}
    An isothermally parametrized surface
    $x=(x^1,x^2,x^3)$ is minimal if and only if all three coordinate
    functions are harmonic:
    \[
        \Delta x^i=x^i_{uu}+x^i_{vv}=0,
        \qquad i=1,2,3.
    \]
\end{corollary}

\begin{example}[Catenoid]
    The parametrization
    \[
        x(u,v)=(a\cosh v\cos u,a\cosh v\sin u,av)
    \]
    satisfies $E=G=a^2\cosh^2v$, $F=0$, and
    $x_{uu}+x_{vv}=0$.  Hence the catenoid is minimal.  Up to pieces of a
    plane, it is the only minimal surface of revolution.
\end{example}

\begin{example}[Helicoid]
    The parametrization
    \[
        x(u,v)=(a\sinh v\cos u,a\sinh v\sin u,au)
    \]
    is isothermal and harmonic, hence minimal.  Apart from the plane, the
    helicoid is the only ruled minimal surface.
\end{example}

\begin{example}[Enneper surface]
    Enneper's minimal surface is
    \[
        x(u,v)=\left(
        u-\frac{u^3}{3}+uv^2,
        v-\frac{v^3}{3}+vu^2,
        u^2-v^2\right).
    \]
    It is a regular minimal parametrized surface with self-intersections.
\end{example}

Let $z=u+iv$ and, for a parametrized surface
$x=(x^1,x^2,x^3)$, set
\[
    \varphi_j=(x^j)_u-i(x^j)_v,
    \qquad j=1,2,3.
\]
Then
\[
    \varphi_1^2+\varphi_2^2+\varphi_3^2=E-G-2iF.
\]
Consequently, $x$ is isothermal if and only if this sum vanishes; under this
condition, $x$ is minimal if and only if the functions $\varphi_j$ are
holomorphic.  This is the starting point for the complex-analytic theory of
minimal surfaces.

\begin{theorem}[Osserman]
    Let $S\subset\mathbb{R}^3$ be a regular minimal surface which is closed as
    a subset of $\mathbb{R}^3$.  If $S$ is not a plane, then the image of its
    Gauss map is dense in $\mathbb{S}^2$.
\end{theorem}




\chapter{The Intrinsic Geometry of Surfaces}

\section{Introduction}

The intrinsic geometry of a surface is the part of its geometry which can be
read from the first fundamental form alone.  Thus notions such as length, angle,
area, covariant derivative, geodesic curvature, and Gaussian curvature are
intrinsic.  The main point of this chapter is that many quantities which are
defined by using the immersion in $\mathbb{R}^3$ actually depend only on the
metric of the surface.

The central results are Gauss' Theorema Egregium, the theory of geodesics and
parallel transport, and the Gauss-Bonnet theorem.  Together they show how local
curvature controls global topological information.

\section{Isometries; Conformal Maps}

\begin{definition}
    Let $S,\bar S$ be regular surfaces.  A diffeomorphism
    $\varphi:S\rightarrow \bar S$ is an \textbf{isometry} if, for every
    $p\in S$ and every $v,w\in T_pS$,
    \[
        \langle v,w\rangle_p
        =
        \langle d\varphi_p(v),d\varphi_p(w)\rangle_{\varphi(p)} .
    \]
    If this condition holds only in a neighborhood of each point, $\varphi$ is
    called a \textbf{local isometry}.
\end{definition}

\begin{proposition}
    Let $x:U\rightarrow S$ and $\bar x:U\rightarrow \bar S$ be parametrizations
    with the same parameter domain.  The map
    $\varphi=\bar x\circ x^{-1}$ is a local isometry if and only if the
    coefficients of the first fundamental forms agree:
    \[
        E=\bar E,\qquad F=\bar F,\qquad G=\bar G .
    \]
\end{proposition}

\begin{remark}
    Isometries preserve lengths of curves, angles between tangent vectors, and
    areas of regions.  They need not preserve quantities depending on the second
    fundamental form.  For example, a plane and a circular cylinder are locally
    isometric although their normal curvatures are different.
\end{remark}

\begin{definition}
    A diffeomorphism $\varphi:S\rightarrow\bar S$ is \textbf{conformal} if
    there is a positive differentiable function $\lambda:S\rightarrow\mathbb{R}$
    such that
    \[
        \langle d\varphi_p(v),d\varphi_p(w)\rangle_{\varphi(p)}
        =
        \lambda(p)^2\langle v,w\rangle_p
    \]
    for all $p\in S$ and all $v,w\in T_pS$.
\end{definition}

\begin{proposition}
    In local coordinates, $\varphi=\bar x\circ x^{-1}$ is conformal if and only
    if
    \[
        \bar E=\lambda^2E,\qquad
        \bar F=\lambda^2F,\qquad
        \bar G=\lambda^2G
    \]
    for some positive function $\lambda$.
\end{proposition}

\begin{remark}
    A conformal map preserves angles, but in general it changes lengths and
    areas by scale factors.  Isometries are the special conformal maps with
    $\lambda=1$.
\end{remark}

\section{The Gauss Theorem and the Equations of Compatibility}

Let $x(u,v)$ be a parametrization of an oriented regular surface and let
$N$ be its unit normal.  Write
\[
    E=\langle x_u,x_u\rangle,\qquad
    F=\langle x_u,x_v\rangle,\qquad
    G=\langle x_v,x_v\rangle
\]
and
\[
    e=\langle x_{uu},N\rangle,\qquad
    f=\langle x_{uv},N\rangle,\qquad
    g=\langle x_{vv},N\rangle .
\]

\begin{definition}
    The \textbf{Christoffel symbols} $\Gamma_{ij}^k$ are defined by the
    tangential parts of the second derivatives of $x$:
    \[
        \begin{aligned}
            x_{uu} & = \Gamma_{11}^1x_u+\Gamma_{11}^2x_v+eN,  \\
            x_{uv} & = \Gamma_{12}^1x_u+\Gamma_{12}^2x_v+fN,  \\
            x_{vv} & = \Gamma_{22}^1x_u+\Gamma_{22}^2x_v+gN .
        \end{aligned}
    \]
\end{definition}

\begin{remark}
    The Christoffel symbols are determined only by $E,F,G$ and their first
    derivatives.  They are therefore intrinsic coefficients of the metric.
\end{remark}

\begin{proposition}
    The Weingarten equations may be written as
    \[
        \begin{aligned}
            N_u & = a_{11}x_u+a_{21}x_v, \\
            N_v & = a_{12}x_u+a_{22}x_v,
        \end{aligned}
    \]
    where the coefficients are determined by
    \[
        \begin{pmatrix}
            a_{11} & a_{12} \\
            a_{21} & a_{22}
        \end{pmatrix}
        =
        -\begin{pmatrix}
            E & F \\
            F & G
        \end{pmatrix}^{-1}
        \begin{pmatrix}
            e & f \\
            f & g
        \end{pmatrix}.
    \]
\end{proposition}

\begin{theorem}[Gauss Equation]
    The coefficients of the first and second fundamental forms satisfy
    \[
        K=\frac{eg-f^2}{EG-F^2}.
    \]
    Moreover, the same function $K$ can be expressed entirely in terms of
    $E,F,G$ and their first and second derivatives.
\end{theorem}

\begin{theorem}[Theorema Egregium]
    The Gaussian curvature of a regular surface is invariant under local
    isometries.  Equivalently, $K$ is determined by the first fundamental form.
\end{theorem}

\begin{remark}
    Although the definition of $K$ uses the Gauss map and the second fundamental
    form, Gauss' theorem says that $K$ is an intrinsic quantity.  This explains
    why a plane and a cylinder are locally isometric: both have $K=0$.
\end{remark}

\begin{theorem}[Compatibility Equations]
    The Gauss equation together with the Mainardi-Codazzi equations gives the
    compatibility conditions for the six functions
    \[
        E,\ F,\ G,\ e,\ f,\ g
    \]
    to arise as the first and second fundamental forms of a surface.
\end{theorem}

\begin{remark}
    The fundamental theorem of local surface theory says that, under the
    compatibility equations and the positivity condition $EG-F^2>0$, these
    data determine a surface locally up to a rigid motion.
\end{remark}

\section{Parallel Transport. Geodesics}

\begin{definition}
    Let $w(t)$ be a differentiable vector field along a curve
    $\alpha:I\rightarrow S$.  The \textbf{covariant derivative} of $w$ along
    $\alpha$ is
    \[
        \frac{Dw}{dt}
        =
        \left(\frac{dw}{dt}\right)^T,
    \]
    the tangential component of the ordinary derivative in $\mathbb{R}^3$.
\end{definition}

\begin{remark}
    If $\alpha(t)=x(u(t),v(t))$ and
    \[
        w(t)=a(t)x_u+b(t)x_v,
    \]
    then
    \[
        \frac{Dw}{dt}
        =
        A x_u+B x_v,
    \]
    where
    \[
        \begin{aligned}
            A & = a'
            +\Gamma_{11}^1au'
            +\Gamma_{12}^1(av'+bu')
            +\Gamma_{22}^1bv', \\
            B & = b'
            +\Gamma_{11}^2au'
            +\Gamma_{12}^2(av'+bu')
            +\Gamma_{22}^2bv'.
        \end{aligned}
    \]
    Hence covariant differentiation depends only on the first fundamental form.
\end{remark}

\begin{definition}
    A vector field $w$ along $\alpha$ is \textbf{parallel} if
    \[
        \frac{Dw}{dt}=0.
    \]
\end{definition}

\begin{proposition}
    Given $w(t_0)\in T_{\alpha(t_0)}S$, there exists a unique parallel vector
    field along $\alpha$ with this initial value.
\end{proposition}

\begin{proposition}
    Parallel transport preserves inner products.  In particular, if $v$ and $w$
    are parallel along $\alpha$, then
    \[
        \frac{d}{dt}\langle v,w\rangle=0.
    \]
\end{proposition}

\begin{definition}
    A parametrized curve $\alpha:I\rightarrow S$ is a \textbf{geodesic} if its
    tangent field is parallel:
    \[
        \frac{D\alpha'}{dt}=0.
    \]
    If $\alpha$ is parametrized by arc length, this is equivalent to saying that
    its acceleration is normal to the surface.
\end{definition}

\begin{remark}
    In local coordinates $\alpha(t)=x(u(t),v(t))$, the geodesic equations are
    \[
        \begin{aligned}
            u''+\Gamma_{11}^1(u')^2
            +2\Gamma_{12}^1u'v'
            +\Gamma_{22}^1(v')^2 & =0,  \\
            v''+\Gamma_{11}^2(u')^2
            +2\Gamma_{12}^2u'v'
            +\Gamma_{22}^2(v')^2 & =0 .
        \end{aligned}
    \]
\end{remark}

\begin{definition}
    Let $\alpha(s)$ be a regular curve on an oriented surface, parametrized by
    arc length.  Put
    \[
        t=\alpha'(s),\qquad n=N\times t .
    \]
    The \textbf{geodesic curvature} of $\alpha$ is
    \[
        k_g(s)=\left\langle \frac{Dt}{ds},n\right\rangle .
    \]
\end{definition}

\begin{proposition}
    A curve parametrized by arc length is a geodesic if and only if
    $k_g=0$.
\end{proposition}

\begin{example}
    The straight lines on a plane and the great circles on a sphere are
    geodesics.  On a circular cylinder, the geodesics are the images of straight
    lines under the local isometry from the plane to the cylinder.
\end{example}

\begin{example}[Clairaut Relation]
    Let a surface of revolution be parametrized by
    \[
        x(u,v)=(f(v)\cos u,f(v)\sin u,g(v)).
    \]
    If $\alpha$ is a geodesic and $\theta$ is the angle between $\alpha'$ and a
    meridian, then
    \[
        f(v)\sin\theta=\text{constant}.
    \]
\end{example}

\section{The Gauss-Bonnet Theorem and Its Applications}

\begin{definition}
    A \textbf{regular region} $R\subset S$ is a compact connected region whose
    boundary is a finite union of piecewise regular simple closed curves.
\end{definition}

\begin{definition}
    Let $\alpha$ be a positively oriented piecewise regular curve and let
    $\theta_i$ be the exterior angle at its vertices.  The sign of $\theta_i$ is
    chosen according to the orientation of the surface.
\end{definition}

\begin{theorem}[Local Gauss-Bonnet Theorem]
    Let $R$ be a simple region contained in a coordinate neighborhood of an
    oriented surface.  Suppose that $\partial R$ is a positively oriented
    piecewise regular curve, parametrized by arc length on each smooth arc.
    Then
    \[
        \sum_i\int_{\alpha_i}k_g\,ds+\sum_i\theta_i+\iint_R K\,d\sigma
        =
        2\pi .
    \]
\end{theorem}

\begin{corollary}
    For a geodesic triangle with interior angles $\varphi_1,\varphi_2,\varphi_3$,
    \[
        \varphi_1+\varphi_2+\varphi_3-\pi
        =
        \iint_R K\,d\sigma .
    \]
\end{corollary}

\begin{remark}
    On a surface of positive curvature, small geodesic triangles have angle sum
    greater than $\pi$; on a surface of negative curvature, the angle sum is
    less than $\pi$.
\end{remark}

\begin{definition}
    A \textbf{triangulation} of a regular region $R$ is a finite decomposition
    of $R$ into triangles whose intersections are either empty, a common vertex,
    or a common edge.
\end{definition}

\begin{definition}
    If a triangulation has $F$ faces, $E$ edges, and $V$ vertices, the
    \textbf{Euler-Poincare characteristic} is
    \[
        \chi(R)=F-E+V .
    \]
\end{definition}

\begin{theorem}[Global Gauss-Bonnet Theorem]
    Let $R$ be a regular region of an oriented surface and suppose that
    $\partial R$ is piecewise regular and positively oriented.  Then
    \[
        \sum_i\int_{\alpha_i}k_g\,ds+\sum_i\theta_i+\iint_R K\,d\sigma
        =
        2\pi\chi(R).
    \]
\end{theorem}

\begin{corollary}
    If $S$ is a compact oriented surface without boundary, then
    \[
        \iint_S K\,d\sigma=2\pi\chi(S).
    \]
\end{corollary}

\begin{remark}
    The integral of the Gaussian curvature over a compact oriented surface is a
    topological invariant.  This is the main bridge between local differential
    geometry and topology in this course.
\end{remark}

\begin{example}
    For the unit sphere, $K=1$ and $\operatorname{area}(S^2)=4\pi$, hence
    \[
        \iint_{S^2}K\,d\sigma=4\pi=2\pi\chi(S^2).
    \]
    Thus $\chi(S^2)=2$.
\end{example}

\begin{proposition}
    A compact oriented surface with positive Gaussian curvature has positive
    Euler characteristic.  In particular, a compact surface of genus $g$ satisfies
    \[
        \chi=2-2g.
    \]
\end{proposition}

\begin{definition}
    Let $v$ be a vector field with an isolated singularity at $p$.  The
    \textbf{index} of $v$ at $p$ is the winding number of the direction of $v$
    along a small positively oriented curve surrounding $p$.
\end{definition}

\begin{theorem}[Poincare-Hopf Formula for Surfaces]
    If $S$ is compact and oriented and $v$ has only isolated singularities, then
    \[
        \sum_p I_p(v)=\chi(S).
    \]
\end{theorem}

\section{The Exponential Map. Geodesic Polar Coordinates}

\begin{definition}
    Let $p\in S$ and $v\in T_pS$.  Let $\gamma_v$ be the geodesic with
    \[
        \gamma_v(0)=p,\qquad \gamma_v'(0)=v.
    \]
    When $\gamma_v(1)$ is defined, the \textbf{exponential map} at $p$ is
    \[
        \exp_p(v)=\gamma_v(1).
    \]
\end{definition}

\begin{remark}
    For small $v$, the exponential map is defined and differentiable.  Moreover,
    \[
        d(\exp_p)_0=\id_{T_pS}.
    \]
    Hence $\exp_p$ is a diffeomorphism from a neighborhood of $0\in T_pS$ onto
    a neighborhood of $p$ in $S$.
\end{remark}

\begin{definition}
    A neighborhood obtained as the image of a disk around $0\in T_pS$ under
    $\exp_p$ is called a \textbf{normal neighborhood} of $p$.
\end{definition}

\begin{proposition}
    In a sufficiently small normal neighborhood, every point $q$ is joined to
    $p$ by a unique geodesic segment, namely the radial geodesic
    \[
        t\longmapsto \exp_p(tv),\qquad 0\leq t\leq 1,
    \]
    where $\exp_p(v)=q$.
\end{proposition}

\begin{definition}
    Choose polar coordinates $(\rho,\theta)$ in $T_pS$.  The map
    \[
        x(\rho,\theta)=\exp_p(\rho e_\theta)
    \]
    gives \textbf{geodesic polar coordinates} on a punctured normal
    neighborhood.
\end{definition}

\begin{proposition}
    In geodesic polar coordinates the first fundamental form has the form
    \[
        d\rho^2+G(\rho,\theta)d\theta^2 .
    \]
    Thus the radial curves $\theta=\text{constant}$ are unit speed geodesics
    and are orthogonal to the geodesic circles $\rho=\text{constant}$.
\end{proposition}

\begin{remark}
    The behavior of $G(\rho,\theta)$ near $\rho=0$ measures the spreading of
    geodesics issuing from $p$.  Curvature controls this spreading: positive
    curvature makes nearby geodesics converge faster, while negative curvature
    makes them separate faster.
\end{remark}

\begin{theorem}[Gauss Lemma]
    In a normal neighborhood of $p$, the radial direction is orthogonal to the
    directions tangent to the geodesic spheres centered at $p$.
\end{theorem}

\begin{proposition}
    Let $q$ be sufficiently close to $p$.  Among all piecewise regular curves
    joining $p$ to $q$ inside a normal neighborhood, the radial geodesic segment
    has minimal length.
\end{proposition}

\begin{definition}
    The \textbf{geodesic distance} between two points is
    \[
        d(p,q)=\inf L(\alpha),
    \]
    where the infimum is taken over all piecewise regular curves joining $p$ to
    $q$.
\end{definition}

\begin{remark}
    Locally, the geodesic distance agrees with the length of the unique radial
    geodesic in a normal neighborhood.
\end{remark}

\section{Further Properties of Geodesics; Convex Neighborhoods}

\begin{theorem}[Existence and Uniqueness of Geodesics]
    Given $p\in S$ and $v\in T_pS$, there exists a unique geodesic
    $\gamma(t)$ defined for $t$ in some interval around $0$ such that
    \[
        \gamma(0)=p,\qquad \gamma'(0)=v.
    \]
    Moreover, $\gamma$ depends differentiably on the initial data $(p,v)$.
\end{theorem}

\begin{remark}
    This follows from the local geodesic equations, which form a system of
    ordinary differential equations with differentiable coefficients.
\end{remark}

\begin{definition}
    A neighborhood $W\subset S$ is \textbf{convex} if any two points of $W$ can
    be joined by a unique geodesic segment contained in $W$, and this segment
    depends differentiably on the two endpoints.
\end{definition}

\begin{proposition}
    Every point of a regular surface has a convex neighborhood.
\end{proposition}

\begin{proposition}
    If $\alpha:[a,b]\rightarrow S$ is a minimizing curve parametrized by arc
    length, then $\alpha$ is a geodesic.
\end{proposition}

\begin{remark}
    Thus geodesics are local candidates for shortest paths.  The converse is
    only local: a geodesic segment need not remain minimizing after it is
    extended far enough.
\end{remark}

\begin{definition}
    A geodesic $\gamma$ is \textbf{closed} if it is periodic as a parametrized
    curve.  It is \textbf{simple closed} if its trace is a simple closed curve.
\end{definition}

\begin{example}
    Great circles on the sphere are closed geodesics.  On a circular cylinder,
    the closed geodesics are the parallels obtained from horizontal lines in
    the plane.  Helices are also geodesics, but they are not closed.
\end{example}

\begin{remark}
    Convex neighborhoods are used to compare local and global behavior.  They
    allow one to replace short pieces of arbitrary curves by unique minimizing
    geodesic segments.
\end{remark}

\section*{Appendix: Proofs of the Fundamental Theorems of the Local Theory of Curves and Surfaces}
\addcontentsline{toc}{section}{Appendix: Proofs of the Fundamental Theorems of the Local Theory of Curves and Surfaces}

\begin{remark}
    The fundamental theorem for curves is proved by solving the Frenet system.
    Given functions $\kappa>0$ and $\tau$, one first constructs an orthonormal
    moving frame satisfying the Frenet equations and then obtains the curve by
    integrating its tangent vector.  Uniqueness follows from uniqueness for
    ordinary differential equations and from the fact that two initial
    orthonormal frames differ by a rigid motion.
\end{remark}

\begin{remark}
    The fundamental theorem for surfaces is analogous but technically richer.
    The prescribed functions $E,F,G,e,f,g$ must satisfy the Gauss and
    Mainardi-Codazzi equations.  These are precisely the integrability
    conditions for the first-order system which reconstructs the frame
    $x_u,x_v,N$.  Once the frame is recovered, the surface is obtained by
    integrating $x_u$ and $x_v$.
\end{remark}

\begin{theorem}[Fundamental Theorem of Local Surface Theory]
    Let $E,F,G,e,f,g$ be differentiable functions on a domain
    $U\subset\mathbb{R}^2$ with
    \[
        E>0,\qquad EG-F^2>0.
    \]
    If the Gauss and Mainardi-Codazzi equations hold, then for each point of
    $U$ there is a neighborhood on which these functions are the coefficients
    of the first and second fundamental forms of a regular surface.  The surface
    is unique up to a rigid motion of $\mathbb{R}^3$.
\end{theorem}

\chapter{Global Differential Geometry}

\section{Introduction}
\section{The Rigidity of the Sphere}
\section{Complete Surfaces. Theorem of Hopf-Rinow}
\section{First and Second Variations of Arc Length; Bonnet's Theorem}
\section{Jacobi Fields and Conjugate Points}
\section{Covering Spaces; The Theorems of Hadamard}
\section{Global Theorems for Curves: The Fary-Milnor Theorem}
\section{Surfaces of Zero Gaussian Curvature}
\section{Jacobi's Theorems}
\section{Abstract Surfaces; Further Generalizations}
\section{Hilbert's Theorem}
\section*{Appendix: Point-Set Topology of Euclidean Spaces}
\addcontentsline{toc}{section}{Appendix: Point-Set Topology of Euclidean Spaces}

\backmatter


\end{document}
